\documentclass[a4paper,11pt]{article}
\usepackage[utf8]{inputenc}
\usepackage{amsmath}
\usepackage{amssymb}
\usepackage{geometry}
\geometry{margin=2.5cm} % Pour avoir de plus belles marges

\title{\textbf{Analyse IV - Summary}}
\author{Liam Vuilleumier}
\date{2026}

\begin{document}
\maketitle

\section{Information}
A concise and rigorous summary transcribing the most important concepts and theorems from the Analyse IV lecture at EPFL.

\section{Holomorphic and Analytic Functions}
\subsection{Definitions}
Let $\Omega \subseteq \mathbb{C}$ be an open set. A complex function $f: \Omega \to \mathbb{C}$ is said to be \textbf{complex differentiable} at a point $z_0 \in \Omega$ if the following limit exists:
\begin{align*}
    f'(z_0) = \lim_{z \to z_0} \frac{f(z) - f(z_0)}{z - z_0}
\end{align*}

If $f$ is complex differentiable at \textit{every} point of $\Omega$, we say $f$ is \textbf{holomorphic} on $\Omega$. 
\\ If $f$ is holomorphic on the entire complex plane ($\Omega = \mathbb{C}$), we call $f$ an \textbf{entire function} (e.g., polynomials, $e^z$, $\sin(z)$).

\section{The Cauchy-Riemann Equations}
Let $z = x + iy \in \mathbb{C}$. A complex function $f(z)$ can be separated into its real and imaginary parts:
\begin{align*}
    f(x+iy) = u(x,y) + i v(x,y)
\end{align*}
where $u, v : \mathbb{R}^2 \to \mathbb{R}$.

\subsection{Necessary and Sufficient Conditions}
\textbf{Theorem:} Suppose $u(x,y)$ and $v(x,y)$ have continuous first-order partial derivatives (i.e., they are of class $C^1$) on an open set $\Omega$. The function $f = u + iv$ is holomorphic on $\Omega$ if and only if it satisfies the \textbf{Cauchy-Riemann equations}:
\begin{align*}
    \frac{\partial u}{\partial x} = \frac{\partial v}{\partial y} \quad \text{and} \quad \frac{\partial u}{\partial y} = -\frac{\partial v}{\partial x}
\end{align*}

When these equations hold, the complex derivative can be computed as:
\begin{align*}
    f'(z) = \frac{\partial u}{\partial x} + i \frac{\partial v}{\partial x} = \frac{\partial v}{\partial y} - i \frac{\partial u}{\partial y}
\end{align*}
\subsection{The Complex Chain Rule}
The chain rule for complex functions works exactly like the chain rule in real calculus. This is a vital tool for proving that a composite function is holomorphic without having to check the Cauchy-Riemann equations directly.

\textbf{The Formula:} \\
Let $g: \Omega \to U$ be holomorphic on $\Omega \subseteq \mathbb{C}$, and let $f: U \to \mathbb{C}$ be holomorphic on $U$. Then the composite function $h(z) = (f \circ g)(z) = f(g(z))$ is holomorphic on $\Omega$, and its derivative is:
\begin{align*}
    h'(z) = f'(g(z)) \cdot g'(z)
\end{align*}

\subsubsection{How to Use the Chain Rule}
To compute the derivative of a complex composition, follow these steps:
\begin{enumerate}
    \item \textbf{Identify the components:} Break the function into an "outer" function $f(w)$ and an "inner" function $g(z)$.
    \item \textbf{Differentiate the outer function:} Compute $f'(w)$ with respect to its variable $w$.
    \item \textbf{Differentiate the inner function:} Compute $g'(z)$ with respect to $z$.
    \item \textbf{Substitute and Multiply:} Multiply the results together, replacing the variable $w$ in $f'(w)$ with the original inner function $g(z)$.
\end{enumerate}

\textbf{Example: Differentiating $h(z) = e^{z^2}$}
\begin{itemize}
    \item Let the outer function be $f(w) = e^w \implies f'(w) = e^w$.
    \item Let the inner function be $g(z) = z^2 \implies g'(z) = 2z$.
    \item Applying the formula:
    \begin{align*}
        h'(z) = e^{g(z)} \cdot (2z) = 2z e^{z^2}
    \end{align*}
\end{itemize}

\subsubsection{Example: Logarithmic Composition}
Consider $f(z) = \log(1+z^2)$ using the principal branch of the logarithm.

\textbf{1. Domain and Holomorphicity:} \\
The principal branch $\log(w)$ is holomorphic for $w \in \mathbb{C} \setminus (-\infty, 0]$. Thus, $f(z)$ is holomorphic where $1+z^2$ is not a non-positive real number. This occurs when $z \neq iy$ for $|y| \geq 1$. The largest open subset of holomorphicity is:
\begin{equation*}
    \Omega = \mathbb{C} \setminus \{iy : y \in \mathbb{R}, |y| \geq 1\}
\end{equation*}

\textbf{2. Derivative via Chain Rule:} \\
Let $w = g(z) = 1+z^2$. Then $f(z) = \log(w)$.
\begin{align*}
    f'(z) &= \frac{d}{dw}(\log w) \cdot \frac{d}{dz}(1+z^2) \\
    &= \frac{1}{1+z^2} \cdot 2z = \frac{2z}{1+z^2}
\end{align*}

\textbf{Note on Holomorphicity:} \\
If $g$ is holomorphic on $\Omega$ and $f$ is holomorphic on $g(\Omega)$, then $f \circ g$ is guaranteed to be holomorphic on $\Omega$. This allows us to conclude that functions like $\sin(e^z)$ or $\cos(z^3)$ are \textbf{entire functions} because they are compositions of entire functions.
\subsection{Harmonic Conjugates}
Because holomorphic functions are infinitely differentiable, $u$ and $v$ have continuous second-order derivatives. By applying the Cauchy-Riemann equations and Schwarz's theorem (symmetry of second derivatives), we get:
\begin{align*}
    \Delta u = \frac{\partial^2 u}{\partial x^2} + \frac{\partial^2 u}{\partial y^2} = 0 \quad \text{and} \quad \Delta v = \frac{\partial^2 v}{\partial x^2} + \frac{\partial^2 v}{\partial y^2} = 0
\end{align*}
Therefore, the real and imaginary parts of a holomorphic function are \textbf{harmonic functions}. We say that $v$ is the harmonic conjugate of $u$.

\section{The Complex Logarithm}
\subsection{Principal Argument}
For any $z \neq 0$, the principal argument $\text{Arg}(z)$ is the unique angle in the interval $(-\pi, \pi]$. It is computed using the arctangent function depending on the quadrant of $(x,y)$. 
\\ \textit{Crucial limitation:} $\text{Arg}(z)$ is discontinuous across the negative real axis. \\
this is how we compute it : 
\[
\text{Arg}(z) =
\begin{cases}
\arctan\left(\dfrac{y}{x}\right) & \text{if } x > 0 \\[6pt]
\arctan\left(\dfrac{y}{x}\right) + \pi & \text{if } x < 0 \text{ and } y \ge 0 \\[6pt]
\arctan\left(\dfrac{y}{x}\right) - \pi & \text{if } x < 0 \text{ and } y < 0 \\[6pt]
\dfrac{\pi}{2} & \text{if } x = 0 \text{ and } y > 0 \\[6pt]
-\dfrac{\pi}{2} & \text{if } x = 0 \text{ and } y < 0
\end{cases}
\]
\subsection{Principal Branch of the Logarithm}
We define the \textbf{principal branch} of the complex logarithm, denoted $\text{Log}(z)$, as:
\begin{align*}
    \text{Log}(z) = \ln|z| + i \text{Arg}(z) = \ln\left(\sqrt{x^2 + y^2}\right) + i \text{Arg}(x+iy)
\end{align*}

\textbf{Rigorous Properties:}
\begin{itemize}
    \item \textbf{Domain:} $\text{Log}(z)$ is a well-defined, single-valued function on the cut plane $\mathbb{C} \setminus (-\infty, 0]$. We must remove zero and the negative real axis (the branch cut) to maintain continuity.
    \item \textbf{Holomorphy:} On this cut plane, $\text{Log}(z)$ is holomorphic.
    \item \textbf{Derivative:} Its complex derivative is exactly what we expect from real analysis: 
    \begin{align*}
        \frac{d}{dz}\text{Log}(z) = \frac{1}{z} \quad \text{for } z \in \mathbb{C} \setminus (-\infty, 0]
    \end{align*}
\end{itemize}
\section{Lecture 3}
\section{Complex Integration}

\textbf{Definition (Line Integral)} \\
Let $\gamma : [a,b] \rightarrow \mathbb{C}$ be a regular (smooth) curve that parametrizes the path $\Gamma \subseteq \mathbb{C}$. Let $f: \Gamma \rightarrow \mathbb{C}$ be a continuous function. We define the line integral of $f$ along $\Gamma$ by:
\begin{equation*}
    \int_\Gamma f(z) \, dz = \int_a^b f(\gamma(t)) \gamma'(t) \, dt
\end{equation*}

\textbf{Example} \\
Let $\Gamma$ be the unit semi-circle centered at $z=0$, parametrized counterclockwise: $\gamma(t) = e^{it}$ for $t \in [0, \pi]$. Then $dz = i e^{it} dt$.

\subsection{Cauchy's Theorem}
Let $U \subseteq \mathbb{C}$ be an open and simply connected set, and let $\gamma$ be a closed and piecewise regular curve inside $U$. If $f: U \rightarrow \mathbb{C}$ is holomorphic, then:
\begin{equation*}
    \oint_\Gamma f(z) \, dz = 0
\end{equation*}

\begin{itemize}
    \item \textbf{Simply connected:} The set consists of "one piece" and has no holes.
    \item \textbf{Piecewise regular:} $\Gamma$ consists of finitely many smooth curves joined end-to-end (e.g., the boundary of a square).
\end{itemize}

\subsubsection{Example}
Let $\Gamma$ be the unit circle and $f(z) = z^2$. Since $f$ is holomorphic on the simply connected set $U = \mathbb{C}$, by Cauchy's Theorem:
\begin{equation*}
    \oint_\Gamma z^2 \, dz = 0
\end{equation*}

\subsection{The Extension of Cauchy's Theorem (Deformation)}
Let $\Gamma_0, \Gamma_1$ be closed, simple, regular curves such that $\Gamma_1 \subseteq \text{int}(\Gamma_0)$. Let $\gamma_0, \gamma_1$ be parametrizations of $\Gamma_0, \Gamma_1$ with the same orientation. If $f$ is holomorphic on the region between $\Gamma_0$ and $\Gamma_1$ (including the boundaries), then:
\begin{equation*}
    \oint_{\Gamma_0} f(z) \, dz = \oint_{\Gamma_1} f(z) \, dz
\end{equation*}

\subsection{Cauchy's Integral Formula}
Suppose $D \subseteq \mathbb{C}$ is open and $f: D \rightarrow \mathbb{C}$ is holomorphic. If $z_0 \in D$, then:
\begin{equation*}
    f(z_0) = \frac{1}{2\pi i} \oint_\gamma \frac{f(z)}{z - z_0} \, dz
\end{equation*}
for every simple, closed, counterclockwise-oriented curve $\gamma$ in $D$ whose interior contains $z_0$.
\section{Lecture 4}
\section{Tricks to Compute Integrals Easily}

\begin{itemize}
    \item \textbf{Cauchy's Theorem:} If $f$ is holomorphic on a \textbf{simply connected} set $D$, then for any closed loop $\gamma$ in $D$:
    \begin{align*}
        \int_{\gamma} f(z) \, dz = 0
    \end{align*}
    
    \item \textbf{Deformation Theorem:} We can replace an integral over a "complicated" curve $\Gamma_0$ by an integral over a "simpler" curve $\Gamma_1$ (like a small circle) as long as $f$ is holomorphic in the region between $\Gamma_0$ and $\Gamma_1$.
    
    \item \textbf{Path Independence:} If $\Gamma_0$ and $\Gamma_1$ share the same endpoints and orientation, and $f$ is holomorphic on the region they enclose, then:
    \begin{align*}
        \int_{\Gamma_0} f(z) \, dz = \int_{\Gamma_1} f(z) \, dz
    \end{align*}
    
    \item \textbf{Generalized Cauchy Integral Formula:} Useful for computing integrals of the form $\frac{f(z)}{(z-z_0)^{n+1}}$. If $f$ is holomorphic on and inside $\gamma$:
    \begin{align*}
        \int_\gamma \frac{f(z)}{(z-z_0)^{n+1}} \, dz = \frac{2 \pi i}{n!} f^{(n)}(z_0)
    \end{align*}
    \textit{Example:} To calculate $\int_\gamma \frac{e^z}{z^3} dz$ around $0$, use $n=2$ and $f(z)=e^z$. The result is $\frac{2\pi i}{2!} f''(0) = \pi i$.
\end{itemize}

\section{Taylor Series and Laurent Series}

\subsection{Taylor Series}
\textbf{Real case:} Let $f: \mathbb{R} \rightarrow \mathbb{R}$ be $N+1$ times differentiable. For any $x_0 \in \mathbb{R}$, the Taylor polynomial is $T^N_{x_0} f(x)$. If $f \in \mathcal{C}^{N+1}$, there exists $\theta \in (0,1)$ such that:
\begin{align*}
    f(x) = \sum_{n=0}^N \frac{f^{(n)}(x_0)}{n!}(x-x_0)^n + \frac{f^{(N+1)}(\theta x_0 + (1-\theta)x)}{(N+1)!}(x-x_0)^{N+1}
\end{align*}

\textbf{Complex case:} Let $f: D \rightarrow \mathbb{C}$ be holomorphic. If the disk $B_R(z_0)$ is contained in $D$, then $f$ is analytic and:
\begin{align*}
    f(z) = \sum_{n=0}^\infty \frac{f^{(n)}(z_0)}{n!}(z-z_0)^n \quad \text{for all } z \in B_R(z_0)
\end{align*}
The largest such $R$ is the \textbf{Radius of Convergence}.

\subsection{Laurent Series}
\textbf{Theorem:} Let $f$ be holomorphic on an annulus (punctured disk) $0 < |z-z_0| < R$. Then:
\begin{align*}
    f(z) = \sum_{n = -\infty}^\infty c_n(z-z_0)^n = \dots + \frac{c_{-2}}{(z-z_0)^2} + \frac{c_{-1}}{z-z_0} + c_0 + c_1(z-z_0) + \dots
\end{align*}
where the coefficients are given by:
\begin{align*}
    c_n = \frac{1}{2 \pi i}\int_\gamma \frac{f(w)}{(w-z_0)^{n+1}}dw
\end{align*}

\textbf{Definition:} The series is split into two parts:
\begin{itemize}
    \item \textbf{Singular part:} Terms with negative powers ($n < 0$).
    \item \textbf{Regular part:} Terms with non-negative powers ($n \geq 0$).
\end{itemize}

\section{Classification of Singularities and Residues}

\subsection{Types of Isolated Singularities}
By looking at the singular part of the Laurent series of $f$ at $z_0$:
\begin{enumerate}
    \item \textbf{Removable Singularity:} No negative powers ($c_n = 0$ for all $n < 0$).
    \item \textbf{Pole of order $k$:} The lowest power is $c_{-k} \neq 0$ (finitely many negative terms).
    \item \textbf{Essential Singularity:} Infinitely many negative terms (e.g., $e^{1/z}$ at $z=0$).
\end{enumerate}


\end{document}